\documentclass[11pt,a4paper]{article}
\usepackage[margin=0.78in]{geometry}
\usepackage[utf8]{inputenc}
\usepackage{microtype}
\usepackage{hyperref}
\setlength{\parindent}{0pt}
\setlength{\parskip}{0.55em}
\hypersetup{colorlinks=true,urlcolor=blue}
\begin{document}
\textbf{\large COVER LETTER}

August 31, 2026

The Editor-in-Chief\\
\textit{Revista Portuguesa de Cardiologia}

\textbf{Subject: Submission of an Original Investigation}

Dear Editor-in-Chief,

Please consider our manuscript, \textbf{"Record-Level Evaluation of a Morlet CNN--Conformer for Five-Class ECG Beat Classification,"} as an Original Investigation in \textit{Revista Portuguesa de Cardiologia}.

Automated ECG classification is increasingly relevant to ambulatory monitoring and cardiovascular telehealth, yet random beat-level splitting and aggregate metrics can overstate performance when rare arrhythmias are unevenly distributed across patients. Our study addresses this methodological issue by evaluating a Morlet CNN--Conformer using a strict record-level holdout of complete MIT--BIH recordings.

Across 15{,}573 held-out three-beat sequences, the CNN--Conformer ranked first among four evaluated architectures on aggregate metrics (accuracy 0.60, macro-F1 0.26, weighted-F1 0.68). Importantly, class-specific analysis showed that the apparent gain was concentrated in Normal and ventricular ectopic beats, whereas SVEB and Fusion recognition remained near zero. We believe this failure-aware result is relevant to readers interested in clinically responsible evaluation of AI-enabled ECG systems, because it demonstrates why record separation, class support, and class-specific metrics must accompany headline performance values.

The manuscript is original, has not been published previously, and is not under consideration by another journal. Both authors have read and approved the submitted version. This work received no specific external funding. \textbf{Declarations of interest: none.} Ethics committee approval and informed consent were not required because the study used only publicly available, fully de-identified MIT--BIH data; no participants were prospectively recruited or contacted and no new identifiable human data were collected. The code, experiment configurations, trained models, and prediction arrays underlying the reported analyses are publicly available through a versioned GitHub software release.

Thank you for considering our work.

Sincerely,

\textbf{Alireza Abbaszadeh}, Corresponding Author\\
Department of Computer Engineering, Mashhad Branch, Islamic Azad University\\
Ostad Yousefi St., Ghasem Abad, Mashhad 91871, Iran\\
Email: \href{mailto:alireza.abbaszadeh8558@iau.ir}{alireza.abbaszadeh8558@iau.ir}\\
Telephone: +98 915 206 2041\\
ORCID: \href{https://orcid.org/0009-0007-8253-6042}{0009-0007-8253-6042}

On behalf of Alireza Abbaszadeh and Vahid Torkzadeh
\end{document}
